% =====================================================================
%  Chapter 16 — Look-Ahead: The Five Ways to Cheat
% =====================================================================

\chapter{Look-Ahead: The Five Ways to Cheat}
\label{ch:look-ahead}

A backtest that uses information the investor could not have had is not a
weak backtest. It is a different object from the one it claims to be, and
no amount of statistical care elsewhere repairs it. This chapter enumerates
the five ways it happens in a system like this one, states the mechanism
by which the book prevents each, and identifies the one that the book
cannot fully prevent.

\section{One: the shifted target}

The first is the subtlest and the most common in home-built systems. A
trailing rolling product looks like a forward return, because the index
runs forwards, and a system trained on it learns to predict the past.

The defence is the shift described in Chapter 15, and the test is
arithmetic: the target at row \(t\) must equal the realised return between
\(t\) and \(t+h\), and must be \emph{unavailable} at \(t\). The book
enforces this in the training-set construction of Chapter 17, where the
rows used to fit the model at date \(t\) are restricted to those with
\(s + h \le t\).

\section{Two: the survivorship-adjusted universe}

The second is the universe. A backtest that holds a fund which did not
exist yet, or that had been liquidated, has made an error of the same kind
as the first, in the cross-section rather than in the time series.

The book's defence is structural: the eligible set at date \(t\) is the set
of funds with a finite price on \(t\), and no allocation can be made to any
other fund. The live window of every fund is reported in Appendix A so
that the claim can be checked rather than trusted.

\section{Three: the revised series}

The third is the macro trap of Chapter 12. A macroeconomic series that has
been revised since the formation date carries information the investor did
not have. The exposure is largest for series that are revised
\emph{quietly} --- employment, GDP, and anything derived from them --- and
it is invisible in the results, because the revised series makes the
historical relationship look better than it was.

The book's defence is to use no revised series at all. Where a macro
quantity is wanted it is proxied from the price series, which is never
revised. The cost is a narrower feature set, and Chapter 14 states it.

\section{Four: the whole-sample statistic}

The fourth is the subtlest of the leaks that survive the obvious defences.
Any statistic computed over the whole sample --- a mean, a standard
deviation, a regression coefficient, a correlation --- carries information
about the future into the past, because the future is part of the sample.

The book's defences are specific. Feature standardisation is computed per
date, on the cross-section that exists on that date. The covariance matrix
is estimated from a trailing window only. The model is refitted at every
formation date on the data available at that date. And the parameter
estimates themselves --- the shrinkage intensity, the ensemble weights ---
are computed from realised outcomes only, with the ensemble's horizon-aware
scoring ensuring that a member's score at date \(t\) uses an outcome that
was realised by \(t\) and not later.

That last one is worth a paragraph of its own, because it is the leak that
a Bayesian ensemble is most likely to have. A member with a three-month
horizon produces a forecast at date \(s\) whose outcome is realised at
\(s+3\). If the ensemble scored that member using the outcome on the date
the forecast was \emph{made}, the member would be scored on the future.
The book's scoring rule uses, at date \(t\), only the forecast made at
\(t - h\) for a member with horizon \(h\), whose outcome is realised
exactly at \(t\). A member with a one-month horizon and a member with a
three-month horizon are therefore both scored on outcomes that were
available at the moment of scoring.

\section{Five: the researcher, not the code}

The fifth is not a bug. It is the one the book cannot prevent, and it is
the reason Chapter 4 exists.

A look-ahead leak that is introduced by a person --- by choosing the
feature set after seeing the results, by choosing the start date after
seeing the results, by choosing which configuration to report --- is
invisible to every code-level defence. No amount of careful implementation
will stop it.

The book's defences against it are procedural, and they are the ones
listed in Chapter 4: the trial family is declared before the test and its
size published; the ablations and the sizing variants are in the family
and not only the configurations that might work; the window sensitivity is
reported rather than a single window; and the result is reported whatever
it is, including a null.

\begin{takeaway}[The five leaks]
\begin{enumerate}
\item The unshifted target: the model trained on a trailing product.
\item The survivorship-adjusted universe: holding a fund that did not
exist.
\item The revised macro series: information that was published later.
\item The whole-sample statistic: a mean, a scale or a coefficient
computed over the entire record.
\item The researcher: a choice of feature set, window or configuration
made after seeing the answer.
\end{enumerate}
\end{takeaway}

\section{The one the book cannot fully prevent}

The fifth is not fully preventable and it is worth being honest about how
much of the book's own confidence rests on it.

The book is written by a machine that has read the literature. The
features it uses --- momentum with a skip, reversal, volatility, trend
shape, drawdown state, beta, correlation, illiquidity --- are the
features a reader who has read the same literature would choose. That is
not independence. The book is not discovering the feature set from
scratch; it is re-implementing a small, well-known set of ideas and asking
whether the resulting system, built honestly and evaluated honestly, does
anything.

This is not a fatal objection, for two reasons.

The first is that the literature's features are the ones that survived
replication testing, and Jensen, Kelly and Pedersen's finding is that most
published factors do replicate. Using them is the rational response to
what is known, not a form of cheating.

The second is that the book's claim is not about the feature set at all.
The claim is about the \emph{architecture}: a stated strategic allocation
that the machine cannot touch, a bounded tilt, a posterior that prices its
own uncertainty, a risk governor, a tax engine, and an evaluation that
subtracts the selection effect. A reader who replaced the twelve features
with their own would keep the architecture, and the book's conclusions
about the architecture would survive. The conclusions about the feature
set would not, and the book does not make any that would not survive.

\begin{keytakeaways}
\item The five leaks are the unshifted target, the survivorship-adjusted
universe, the revised macro series, the whole-sample statistic, and the
researcher.
\item The first four are prevented structurally and each defence is stated
where it is made.
\item The fifth cannot be prevented by code; it is prevented procedurally,
by declaring the trial family before the test and reporting the result
whatever it is.
\item The book's own feature set is not independent of the literature, and
the book's claim is about the architecture rather than about the
features.
\end{keytakeaways}
